% Volume VII: Formal Synthesis
% Part XIV. Dynamics, Geometry, and Holonomy
% Chapter 77: Belief-State Manifolds
% Spine: proof-spines/ch077-spine.md

\chapter{Belief-State Manifolds}
\label{ch:ch077}

Part XIV begins by modeling institutional states as points on a \textbf{belief-state manifold} \(\mathcal M\). In that picture, evidence produces local motion, but procedural and legal rules determine which directions are actually available. Not every imaginable revision is institutionally admissible; rather,
\[
\dot\Theta_t\in\mathcal D_{\Theta_t}\subseteq T_{\Theta_t}\mathcal M.
\]
\Definition\ Here \(T_{\Theta_t}\mathcal M\) is the tangent space at the current state \(\Theta_t\), and \(\mathcal D\) is a distribution encoding evidentiary and procedural constraints on motion.

The geometric point is not decorative. Once belief is placed on \(\mathcal M\), one can distinguish small local adjustments from prohibited or normatively dangerous transitions, and one can describe institutional learning as travel through a constrained space rather than as isolated proposition swaps. This chapter therefore provides the formal setting for the later claims about path dependence, curvature, and holonomy: the update problem is always a problem of constrained motion on a socially structured manifold.
